% P05.1 zero-degree literal regularity, 26 September 2026.

\begin{theorem}
\label{thm:frontier-polyhedral-zero-degree-integral-regularity}
Let $\mathbf X$ be a one-dimensional balanced rational polyhedral space
admitting a finite structure, and let
$\gamma\in\mathrm{PAPC}_{\mathbf Z}(\mathbf X)$ be positive at infinity.
Then the following conditions are equivalent:
\begin{enumerate}
\item $\deg(\gamma)=0$;
\item $\mathbf X$ is compact.
\end{enumerate}
When these conditions hold, every function in
$\mathrm{PC}(\mathbf X,\gamma)$ is constant on each connected component.
For every bounded literal piecewise-affine function
$u\in\mathrm{PA}_{\mathbf Z,b}(\mathbf X)$ and every connected component
$C$ of $\mathbf X$, one has
\begin{equation}
 P_\gamma(u)|_C
   =\min_{x\in C}\bigl(\gamma(x)+u(x)\bigr).
 \label{eq:frontier-polyhedral-zero-degree-envelope}
\end{equation}
The right-hand side is a rational constant.  In particular,
$P_\gamma(u)\in\mathrm{PAPC}_{\mathbf Z}(\mathbf X,\gamma)$ and
$(\mathbf X,\gamma)$ has the literal integral-slope regularity property.
\end{theorem}

Thus the zero-degree branch of the one-dimensional literal problem is
regular.  Together with
\Cref{thm:frontier-polyhedral-strict-vertex-integral-nonregularity}, this
leaves only positive-degree graphs for which every positive-mass vertex has
a higher-dimensional local relation cone and the reference slope vector lies
on a proper face of its dual cone.
